Let us assume that we observe a hot Jupiter planet orbiting around a star at an
average distance $d = 5$\,AU. It has been found that the distance of this
system from us is $r = 250$\,pc. What is the minimum diameter, $D$, that a
telescope should have to be able to resolve the two objects (star and planet)?
We assume that the observation is done in the optical part of the
electromagnetic spectrum ($\lambda \sim 500$\,nm), outside the Earth's
atmosphere and that the telescope optics are perfect.
